\documentclass[12pt]{article}
\usepackage[utf8]{inputenc}
\usepackage[T1]{fontenc}
\usepackage{hyperref}
\title{Turkana Language Resources for K-12 Curriculum Development}
\author{Compiled by Claude Sonnet 5 (claude-sonnet-5) for Andr\'as Kornai}
\date{}
\begin{document}
\maketitle

\section{Introduction}
Turkana (Ng'aturkana, ISO 639-3 tuv) is an Eastern Nilotic language spoken mainly in Turkana County, Kenya. Speaker estimates vary by source (about 990,000 in the 2009 Kenyan census figure cited by Wikipedia, higher in regional totals). It is a tone language. Search found a reference grammar, dictionaries, scripture resources and some mother-tongue education research, but no ready-made school textbooks.

\section{Primary Educational Resources}
\begin{itemize}
\item African Storybook: a page on Dr John Ng'asike's work with children in Turkana. It reports that teachers often treat mother-tongue subjects as non-examinable. \url{https://www.africanstorybook.org/use.php?id=25}
\item Ng'asike's research on mother-tongue early childhood education and family literacy storybooks, hosted by the Kenyatta University repository. \url{https://kerd.ku.ac.ke/handle/123456789/803}
\end{itemize}

\section{Linguistic and Reference Materials}
\begin{itemize}
\item Dimmendaal, \emph{The Turkana Language} (1983), the main reference grammar, described as open access. I did not see a direct link. A Cambridge Language Centre summary page is here: \url{https://www.langcen.cam.ac.uk/resources/langt/turkana.html}
\item Cambridge Language Centre, The Turkana language, data page. \url{https://www.langcen.cam.ac.uk/resources/langt/data/tu_101.html}
\item Wikipedia, Turkana language. It lists A. Barrett's English-Turkana (1988) and Turkana-English (1990) dictionaries. \url{https://en.wikipedia.org/wiki/Turkana_language}
\end{itemize}

\section{Religious and Community Resources}
\begin{itemize}
\item Scripture Earth page for Turkana, listing a Bible Society of Kenya audio recording (2002), a Luke film, a Bible app and a YouVersion entry. I did not verify current availability. \url{https://www.scriptureearth.org/00eng.php?iso=tuv}
\end{itemize}

\section{Digital and Audio Resources}
\begin{itemize}
\item Global Recordings Network: Turkana audio materials. \url{https://glorecnet-website.k.globalrecordings.net/en/language/745}
\end{itemize}

\section{Curriculum Development Context}
Turkana is not recorded as a school subject in the survey data, though research shows mother-tongue use in early childhood education is being studied. Barriers noted in the sources include the lack of locally relevant materials, mother-tongue subjects not being examined, and children reluctant to speak Ng'aturkana in front of teachers. A curriculum could build on the Dimmendaal grammar, the Barrett dictionaries and community-collected storybooks.
\end{document}